\renewcommand{\theequation}{D\thechapter.\arabic{equation}}
\label{sec:appendixD}

The effects of using pulse or saturation irradiation data can be demonstrated through an illustrative toy problem.
Instead of using several hundred \glspl{DNP}, only four will be considered: $A, B, C,$ and $D$.
Each of these \glspl{DNP} will have an emission probability $P_n$ of unity and will only be formed via fission, having an equivalent cumulative fission yield $CFY$ and independent fission yield $IFY$ of 0.25.
The only difference between these \glspl{DNP} will be in their half-lives, given as: $\lambda_A=0.001$, $\lambda_B=0.01$, $\lambda_C=0.1$, and $\lambda_D=1$ $s^{-1}$.

The concentration of each nuclide can be calculated using Equation \eqref{eq:cfy-conc}, recreated here:
\begin{equation}
    N_i = \frac{CFY_i}{\lambda_i}.
    \nonumber
\end{equation}
This gives concentrations of $N_A = 250$, $N_B = 25$, $N_C = 2.5$, and $N_D = 0.25$ atoms.

Because all of the \glspl{DNP} come directly from fission, the independent fission yield can be used to calculate the concentration after a single normalized fission event.
This gives concentrations of $N_A = N_B =N_C =N_D=0.25$ atoms.

Next, to complete the comparison, these four \glspl{DNP} can be fit into two \gls{DNP} groups to determine what parameters best fit the delayed neutron count rates produced.
The total delayed neutron yield should be one because one delayed neutron is released per fission event, and the average half-life is calculated to be approximately 192.5 seconds.
Discussion on how these values are calculated is given in Section \ref{sec:gandkparams}.
The results comparing the data of these four \glspl{DNP} to the different group fits are given in Table \ref{tab:pulse-sat-compare}.

\begin{table}[!ht]
    \centering
    \caption{A comparison of the raw \gls{DNP} data to the different irradiation methods and the corresponding yields and average half-lives.}
    \begin{tabular}{ccc}
    \hline
       Method  & Delayed Yield [-] & Average Half-life [$s$] \\
       \hline
        Data & 1.00 & 192.5 \\
        4 Group Saturation & 1.00 & 192.5 \\
        3 Group Saturation & 0.98 & 188.9 \\
        2 Group Saturation & 0.92 & 182.3 \\
        1 Group Saturation & 0.62 & 223.8 \\
        4 Group Pulse & 1.00 & 192.5 \\
        3 Group Pulse & 0.91 & 91.65 \\
        2 Group Pulse & 0.80 & 61.0 \\
        1 Group Pulse & 0.42 & 1.3 \\
        \hline
    \end{tabular}
    \label{tab:pulse-sat-compare}
\end{table}

Because each \gls{DNP} is equally important, the 3 group structure captures 75\% of the total delayed neutrons, while the 2 group structure is only able to capture 50\% of the total delayed neutrons.
The delayed neutron results show that both the pulse and saturation irradiations generally capture the delayed neutrons, provided there are sufficiently many groups available.
As the number of groups decreases, the delayed yield decreases.

The average half-life shows that the saturation irradiation tends to over-estimate the average half-life as the number of groups becomes too few, though until that point it has more accurate values for the yield and average half-life.
The pulse irradiation tends to under-estimate the average half-life noticeably, and the yield decreases faster than the saturation irradiation results.

Although this example does not exactly represent reality, it provides a useful understanding of the effects of both the pulse and saturation irradiations on the \gls{DNP} group parameters.
The results here provide further motivation to investigate combining the pulse and saturation irradiation data in the \gls{DNP} group parameter formulation, as discussed in Section \ref{sec:proposed-0d-flow}.